% !TEX root = ../../../main.tex
% Block card: preprocessing.wavelet_bands -- "Wavelet bands" (Part I, the blocks chapter).
\cardfield{What it does}
Splits a signal into a ladder of layers, each holding movements about half as fast as the
layer above, and passes one chosen layer on down the chain.

\cardfield{When to use it}
When you want to look at one time-scale of activity but do not know in advance where to cut.
It works like sorting coins by size through a stack of sieves: each tray catches one size, and
all the trays together hold every coin you started with. Compared with
\blockref{preprocessing.bandpass}, it keeps a sharp event sharper and in exactly the right
place in time, and every layer lines up sample for sample with the recording. Discovery can
also run a template once per layer (a wavelet band scope).

\cardfield{What goes in, what comes out}
In goes a span of one channel with no missing samples. Out comes one layer, as an ordinary
signal of the same length and rate. The block page draws every layer, with its range of
speeds in hertz, so you can see which one carries your events.

\cardfield{Settings that matter}
\emph{level} picks the layer that goes on: 1 is the fastest (from a quarter to half the
sampling rate), each higher number half as fast, and 0 is the slow remainder below the deepest
layer. \emph{levels} is how many layers to make; leave it at 0 and the block chooses. On a
1\,Hz recording the default level 4 holds movements of roughly 16 to 32 seconds per cycle
(0.031--0.062\,Hz). \emph{wavelet} is the shape used to cut; the default db4 is a good
general choice.

\cardfield{Watch out for}
A span with any missing sample is refused: fill or avoid gaps first. The deepest layers are
distorted near both ends of the span. The edges between layers are approximate: a little of
each layer leaks into its neighbours.
